% Actual class: 01
% Slide decks: L0 Introduction Slides (Week 1); L1 Imaging v1 full
% Exact scope: L0 pp. 1--21; L1 pp. 1--24
% NTULearn supplement: Not used
% Human verified: Yes

\section{第 01 节课：Imaging、Photometry 与 Digital Image Formation}
\label{lecture:SC4061-L01}

\lecturemeta
  {SC4061-L01}
  {2026-08-11}
  {L0 Introduction Slides (Week 1)；L1 Imaging v1 full}
  {L0 pp.~1--21；L1 pp.~1--24}
  {未使用}
  {已确认（2026-08-11）}

\subsection{本讲主线}

本讲说明 digital image（数字图像）如何形成：三维 scene（场景）中的光经过 surface reflection（表面反射）、lens/aperture（透镜/光圈）和 sensor（传感器），再经 sampling（采样）与 quantization（量化）成为矩阵。Computer Vision（计算机视觉）则尝试从最终图像反推出场景信息。

\lecturetopic{SC4061-L01-T01}{Computer Vision 是一个 Inverse Problem}

\keyidea{Image（图像）不是 scene（场景）的直接副本，而是 lighting（光照）、material（材质）、geometry（几何）、optics（光学系统）和 digitization（数字化）共同作用的结果。}

伪装场景说明 local appearance（局部外观）不足以可靠确定 object boundary（物体边界），识别还需要 global shape（整体形状）和 context（上下文）。光照错觉则说明相同 measured intensity（测量强度）在不同环境中可能对应不同的 perceptual interpretation（感知解释）。

\begin{figure}[htbp]
  \centering
  \resizebox{\textwidth}{!}{\begin{tikzpicture}[
  >=Latex,
  stage/.style={draw=noteBlue,rounded corners=2pt,fill=blue!5,
    minimum width=2.7cm,minimum height=1.2cm,align=center},
  flow/.style={->,thick,draw=noteBlue},
  inverse/.style={<-,dashed,very thick,draw=ntuRed}
]
  \node[stage] (scene) {3D scene\\（三维场景）};
  \node[stage,right=0.55cm of scene] (radiance) {lighting \& reflection\\（光照与反射）};
  \node[stage,right=0.55cm of radiance] (optics) {lens \& aperture\\（透镜与光圈）};
  \node[stage,right=0.55cm of optics] (sensor) {sensor exposure\\（传感器曝光）};
  \node[stage,right=0.55cm of sensor] (digital) {sampling + quantization\\digital image（数字图像）};

  \draw[flow] (scene) -- (radiance);
  \draw[flow] (radiance) -- (optics);
  \draw[flow] (optics) -- (sensor);
  \draw[flow] (sensor) -- (digital);

  \draw[inverse]
    ([yshift=-0.25cm]scene.south) to[bend right=22]
    node[below,align=center]{Computer Vision（计算机视觉）：\\从 image evidence（图像证据）反推 scene properties（场景属性）}
    ([yshift=-0.25cm]digital.south);
\end{tikzpicture}
}
  \caption{从 scene（场景）到 digital image（数字图像）的 forward process（正向过程），以及 Computer Vision（计算机视觉）的反向推断目标。}
  \label{fig:sc4061-l01-imaging-pipeline}
\end{figure}

由于多个不同场景可能产生相似图像，从 image（图像）恢复 scene properties（场景属性）通常是 ill-posed inverse problem（不适定逆问题），需要先验知识、上下文或多个观测约束答案。

\lecturetopic{SC4061-L01-T02}{Camera Imaging System}

一个简化的 camera imaging system（相机成像系统）包含 aperture（光圈）、optical system（光学系统）和 image sensor（图像传感器）。

\begin{center}
\small
\begin{tabularx}{\textwidth}{@{}>{\raggedright\arraybackslash}p{2.5cm}>{\raggedright\arraybackslash}p{2.7cm}>{\raggedright\arraybackslash}p{3.1cm}>{\raggedright\arraybackslash}X@{}}
  \toprule
  Human eye（人眼） & Camera（相机） & 核心变量 & 作用 \\
  \midrule
  pupil（瞳孔） & aperture（光圈） & aperture diameter（光圈直径） & 控制进光量，并影响 depth of field（景深） \\
  eye lens（晶状体） & camera lens（相机镜头） & focal length（焦距） & 改变光线方向，在 image plane（像平面）形成聚焦 \\
  retina（视网膜） & sensor（传感器） & exposure（曝光量） & 把到达的光能转换为可数字化的电信号 \\
  \bottomrule
\end{tabularx}
\end{center}

Optical axis（光轴）是系统的中心参考轴。Aperture（光圈）限制允许进入的 ray bundle（光线束），lens（透镜）决定几何映射，sensor plane（传感器平面）记录 image irradiance（图像辐照度）。

\lecturetopic{SC4061-L01-T03}{Thin Lens Model 与 Magnification}

Thin lens model（薄透镜模型）用三个理想化 principal rays（主光线）确定物点的像点：平行于光轴的光线折射后通过远焦点；通过近焦点的光线折射后平行于光轴；通过 optical center（光心）的光线在薄透镜近似下保持直线。

\begin{figure}[htbp]
  \centering
  \resizebox{0.94\textwidth}{!}{\begin{tikzpicture}[>=Latex,font=\small]
  \draw[->,thin] (-6,0) -- (3.3,0) node[right] {optical axis（光轴）};

  % Thin lens and focal points.
  \draw[very thick,draw=noteBlue]
    (0,-2.15) .. controls (-0.28,-1.2) and (-0.28,1.2) .. (0,2.15)
    .. controls (0.28,1.2) and (0.28,-1.2) .. (0,-2.15);
  \fill[noteBlue] (-1.55,0) circle (2pt) node[below] {$F_l$};
  \fill[noteBlue] (1.55,0) circle (2pt) node[above] {$F_r$};
  \fill[noteBlue] (0,0) circle (2pt) node[below left] {$O$};

  % Object and inverted image.
  \draw[->,very thick,draw=ntuRed] (-5,0) -- node[midway,left] {$Y$} (-5,2.4)
    node[above] {$P$};
  \draw[->,very thick,draw=ntuRed] (2.2,0) -- node[midway,right] {$y$} (2.2,-1.03)
    node[below] {$p$};
  \node[draw=gray,rounded corners=1pt,fill=white,inner sep=2pt]
    at (-4.25,-0.58) {$X:\ \odot$};
  \node[draw=gray,rounded corners=1pt,fill=white,inner sep=2pt]
    at (2.72,0.58) {$x:\ \otimes$};

  % Three principal rays.
  \draw[thick,draw=orange!90!black] (-5,2.4) -- (0,2.4) -- (2.2,-1.03);
  \draw[thick,draw=noteBlue] (-5,2.4) -- (0,0) -- (2.2,-1.03);
  \draw[thick,draw=green!45!black] (-5,2.4) -- (0,-1.03) -- (2.2,-1.03);

  % Distance guides.
  \draw[<->] (-5,-1.72) -- node[fill=white,inner sep=1pt] {$u=\hat Z$} (0,-1.72);
  \draw[<->] (0,-1.72) -- node[fill=white,inner sep=1pt] {$v=\hat z$} (2.2,-1.72);
  \draw[<->,draw=gray] (-1.55,-1.35) -- node[fill=white,inner sep=1pt] {$f$} (0,-1.35);
  \draw[<->,draw=gray] (0,-1.35) -- node[fill=white,inner sep=1pt] {$f$} (1.55,-1.35);
\end{tikzpicture}
}
  \caption{Thin lens（薄透镜）的三条 principal rays（主光线）。$Y/y$ 位于图示截面内；$X/x$ 垂直于图示截面，$\odot/\otimes$ 分别表示指向纸外/纸内。}
  \label{fig:sc4061-l01-thin-lens}
\end{figure}

\begin{samepage}
令 $u=\hat Z$ 为 object distance（物距），$v=\hat z$ 为 image distance（像距），$f$ 为 focal length（焦距），则
\[
  \boxed{\frac{1}{u}+\frac{1}{v}=\frac{1}{f}}.
\]
Signed magnification（带符号放大率）为
\[
  \boxed{m=\frac{y}{Y}=\frac{x}{X}=-\frac{v}{u}}.
\]
负号表示 real image（实像）相对物体倒立。讲义图中的 $X/x$ 是垂直于所画截面的另一组横向坐标，与 $Y/y$ 满足相同缩放关系。
\end{samepage}

\textbf{对应例题：}\relatedexercise{SC4061-L01-E01}{Thin Lens 成像与放大率}。

\lecturetopic{SC4061-L01-T04}{Depth of Field 与 Aperture Trade-off}

严格聚焦的 object plane（物平面）会映射为 sensor 上的理想 image plane（像平面）。其他深度的物点在 sensor 上形成 circle of confusion（弥散圆）；当其小于可接受阈值时，视觉上仍被认为清晰。Depth of field（景深）就是满足该清晰标准的物距范围。

对于 off-axis point（离轴点）$P$，通过 optical center（光心）$O$ 的 chief ray（主光线）在 thin-lens approximation（薄透镜近似）下不偏折，因此理想像点 $q$ 必须位于该直线上，而不会落到 optical axis（光轴）上。Sensor plane（传感器平面）与 $q$ 所在的 ideal image plane（理想像平面）不重合时，才形成以 chief ray 与 sensor 交点为中心的 blur circle（模糊圆）。

\begin{figure}[htbp]
  \centering
  \resizebox{0.88\textwidth}{!}{\begin{tikzpicture}[>=Latex,font=\small]
  \draw[->,thin] (-4.8,0) -- (6.7,0) node[right] {optical axis（光轴）};
  \draw[very thick,draw=noteBlue]
    (0,-1.75) .. controls (-0.22,-0.9) and (-0.22,0.9) .. (0,1.75)
    .. controls (0.22,0.9) and (0.22,-0.9) .. (0,-1.75);
  \fill[noteBlue] (0,0) circle (2pt) node[above right] {$O$};

  \draw[dashed,draw=gray] (3.7,-1.85) -- (3.7,1.45);
  \node[above,align=center] at (3.7,1.45) {ideal image\\plane};
  \draw[very thick] (5.0,-2.0) -- (5.0,1.55);
  \node[above] at (5.0,1.55) {sensor plane};

  \fill[ntuRed] (-4,1.0) circle (2.5pt)
    node[above] {off-axis scene point $P$（离轴场景点）};
  \fill[black] (3.7,-0.925) circle (2.2pt)
    node[below right] {$q$};

  % The chief ray passes through the optical center without bending.
  \draw[dashed,thick,draw=gray] (-4,1.0) -- (0,0) -- (5.0,-1.25);

  % Wide aperture ray bundle.
  \draw[thick,draw=ntuRed] (-4,1.0) -- (0,1.1) -- (3.7,-0.925) -- (5.0,-1.636);
  \draw[thick,draw=ntuRed] (-4,1.0) -- (0,-1.1) -- (3.7,-0.925) -- (5.0,-0.864);
  \draw[line width=3.2pt,draw=ntuRed] (5.0,-1.636) -- (5.0,-0.864);

  % Narrow aperture ray bundle.
  \draw[thick,draw=noteBlue] (-4,1.0) -- (0,0.4) -- (3.7,-0.925) -- (5.0,-1.391);
  \draw[thick,draw=noteBlue] (-4,1.0) -- (0,-0.4) -- (3.7,-0.925) -- (5.0,-1.109);
  \draw[line width=3.2pt,draw=noteBlue] (5.07,-1.391) -- (5.07,-1.109);

  \node[text=ntuRed,anchor=west,align=left] at (5.28,-0.70) {larger\\blur circle};
  \node[text=noteBlue,anchor=west,align=left] at (5.28,-1.72) {smaller\\blur circle};

  \draw[very thick,draw=ntuRed] (-1.5,-1.55) -- (-0.8,-1.55)
    node[right,text=black]{wide aperture（大光圈）};
  \draw[very thick,draw=noteBlue] (-1.5,-1.90) -- (-0.8,-1.90)
    node[right,text=black]{narrow aperture（小光圈）};
  \draw[dashed,thick,draw=gray] (-1.5,-2.24) -- (-0.8,-2.24)
    node[right,text=black]{chief ray（主光线）};
\end{tikzpicture}
}
  \caption{离轴场景点 $P$ 聚焦到离轴像点 $q$；sensor plane（传感器平面）偏离理想像平面时，wide aperture（大光圈）产生更大的 blur circle（模糊圆）。}
  \label{fig:sc4061-l01-depth-of-field}
\end{figure}

\[
  \begin{gathered}
    \text{smaller aperture（更小光圈）}
    \Longrightarrow \text{smaller blur circle（更小模糊圆）}\\
    \Longrightarrow \text{larger DoF（更大景深）}
  \end{gathered}
\]
代价是进光量减少，通常需要 longer exposure（更长曝光）或 higher gain（更高增益）；光圈过小时还会受到 diffraction（衍射）限制。Pinhole camera（针孔相机）是 aperture 趋近于零的极端情况，景深大但进光量很小。

\lecturetopic{SC4061-L01-T05}{Photometry：从光照到 Pixel Intensity}

在 Lambertian reflection（朗伯反射）假设下，scene point（场景点）$P$ 的 radiance（辐亮度）可写为
\[
  L(P)=\rho\max\!\left(0,\bm I^{\mathsf T}\bm n\right),
\]
其中 $\rho$ 是 albedo（反照率），$\bm I$ 是 source radiance vector（光源辐亮度向量），$\bm n$ 是 surface normal（表面法向量）。Dot product（点积）描述表面朝向与入射光方向的匹配程度。

\begin{figure}[htbp]
  \centering
  \resizebox{\textwidth}{!}{\begin{tikzpicture}[
  >=Latex,
  stage/.style={draw=noteBlue,rounded corners=2pt,fill=blue!5,
    minimum width=3.1cm,minimum height=1.15cm,align=center},
  flow/.style={->,thick,draw=noteBlue}
]
  \node[stage] (light) at (0,1.05) {light source\\$\bm I$};
  \node[stage] (surface) at (0,-1.05) {surface properties\\$\rho,\bm n$};
  \node[stage] (radiance) at (4.4,0) {scene radiance\\$L(P)=\rho\max(0,\bm I^{\mathsf T}\bm n)$};
  \node[stage] (optics) at (9.0,0) {camera optics\\$d,v,\alpha$};
  \node[stage] (irradiance) at (13.5,0) {sensor irradiance\\$E(p)$};

  \draw[flow] (light.east) -- (radiance.west);
  \draw[flow] (surface.east) -- (radiance.west);
  \draw[flow] (radiance) -- (optics);
  \draw[flow] (optics) -- (irradiance);

  \node[below=0.35cm of optics,align=center]
    {$E(p)=L(P)\dfrac{\pi}{4}\left(\dfrac{d}{v}\right)^2\cos^4\alpha$};
\end{tikzpicture}
}
  \caption{Photometric image formation（光度成像）的两个阶段：表面产生 radiance，光学系统把它转换为 sensor irradiance。}
  \label{fig:sc4061-l01-photometry}
\end{figure}

在讲义采用的理想模型中，sensor point（传感器点）$p$ 的 irradiance（辐照度）为
\[
  E(p)=L(P)\frac{\pi}{4}\left(\frac{d}{v}\right)^2\cos^4\alpha,
\]
其中 $d$ 是 aperture diameter（光圈直径），$v$ 是 image distance（像距），$\alpha$ 是光线相对 optical axis（光轴）的夹角。该式说明进光量近似随 $d^2$ 增加，离轴位置因 $\cos^4\alpha$ 出现 vignetting（暗角）。

更常用的简化表达是
\[
  f(x,y)=i(x,y)r(x,y),
\]
其中 $i$ 为 illumination（光照），$r$ 为 reflectance（反射率）。只观察 $f$ 无法唯一分离 $i$ 与 $r$，这正是 inverse problem（逆问题）的一个具体来源。

\lecturetopic{SC4061-L01-T06}{Sensor、Sampling、Quantization 与 Image Matrix}

Sensor array（传感器阵列）由二维 photo-detectors（光电探测器）组成。每个探测器记录的信号与 exposure（曝光量）相关：
\[
  \text{exposure}=\text{irradiance}\times\text{exposure time}.
\]
单个探测器主要测量强度；Bayer mask（拜耳滤色阵列）让不同位置分别采集 R/G/B（红/绿/蓝）通道，再通过 demosaicing（去马赛克）估计每个 pixel（像素）的完整颜色。

\begin{figure}[htbp]
  \centering
  \resizebox{0.94\textwidth}{!}{\begin{tikzpicture}[>=Latex,font=\small]
  % Sampling panel.
  \begin{scope}
    \draw[->] (0,0) -- (5.4,0) node[right] {$x$};
    \draw[->] (0,0) -- (0,3.6) node[above] {intensity};
    \draw[very thick,draw=orange!90!black,smooth]
      plot coordinates {(0.25,0.55) (0.8,1.25) (1.35,0.9) (1.9,2.35)
                        (2.45,1.55) (3.0,3.05) (3.55,1.25) (4.1,2.15)
                        (4.65,1.65) (5.05,2.65)};
    \foreach \x/\y in {0.5/0.85,1.1/1.05,1.7/1.8,2.3/1.7,2.9/2.8,3.5/1.3,4.1/2.15,4.7/1.75}{
      \draw[->,draw=noteBlue] (\x,0) -- (\x,\y);
      \fill[noteBlue] (\x,\y) circle (1.8pt);
    }
    \node[below=0.55cm,align=center] at (2.7,0)
      {sampling（采样）\\$x=m\Delta x$};
  \end{scope}

  \draw[->,very thick,draw=ntuRed] (5.9,1.8) -- node[above] {$Q(\cdot)$} (7.1,1.8);

  % Quantization panel.
  \begin{scope}[xshift=7.5cm]
    \draw[->] (0,0) -- (5.4,0) node[right] {$m$};
    \draw[->] (0,0) -- (0,3.6) node[above] {level};
    \foreach \y in {0.8,1.6,2.4,3.2}{
      \draw[dotted,draw=gray] (0,\y) -- (5.1,\y);
    }
    \foreach \x/\y in {0.5/0.8,1.1/0.8,1.7/1.6,2.3/1.6,2.9/3.2,3.5/1.6,4.1/2.4,4.7/1.6}{
      \draw[->,draw=noteBlue] (\x,0) -- (\x,\y);
      \fill[noteBlue] (\x,\y) circle (1.8pt);
    }
    \node[below=0.55cm,align=center] at (2.7,0)
      {quantization（量化）\\continuous value $\rightarrow$ finite levels};
  \end{scope}
\end{tikzpicture}
}
  \caption{Sampling（采样）离散化空间位置，quantization（量化）把测量值映射到有限灰度等级。}
  \label{fig:sc4061-l01-sampling-quantization}
\end{figure}

二维数字化过程可统一写成
\[
  \boxed{g[m,n]=Q\!\left(f(m\Delta x,n\Delta y)\right)}.
\]
Sampling interval（采样间隔）过大会丢失空间细节并引发 aliasing（混叠）；quantization levels（量化等级）过少会产生 banding（色带）或明显的强度跳变。$b$-bit quantization（$b$ 位量化）共有 $2^b$ 个等级，例如 8-bit 灰度通常为 $0$--$255$。

Grayscale image（灰度图像）可视为 scalar function（标量函数）$f(x,y)$ 或二维矩阵；RGB image（彩色图像）可视为三个通道组成的三维数组。代码中通常采用
\[
  \boxed{I[\text{row},\text{column}]=I[y,x]},
\]
即 $x$ 对应 column（列），$y$ 对应 row（行）。Python/NumPy 通常从 0 编号，MATLAB 通常从 1 编号。

PNG/TIFF 常用于 lossless representation（无损表示），JPEG 使用 lossy compression（有损压缩），JPEG~2000 使用 wavelet-based compression（基于小波的压缩）。文件大小不能只由 pixel count（像素数）决定，还必须说明 channel count（通道数）、bit depth（位深）和 compression（压缩方式）。

\textbf{对应例题：}\relatedexercise{SC4061-L01-E02}{Sampling、Quantization 与存储量}。

\subsection{一分钟回顾}

三维 scene（场景）先通过 lighting 与 reflection（光照与反射）产生 radiance（辐亮度），再由 lens 和 aperture（透镜和光圈）映射为 sensor irradiance（传感器辐照度），最后经 sampling 与 quantization（采样与量化）成为 image matrix（图像矩阵）。Thin lens equation（薄透镜公式）描述几何成像；aperture（光圈）同时影响亮度和 depth of field（景深）；pixel intensity（像素强度）并不等于物体固有属性。Computer Vision（计算机视觉）则试图沿这条链反向推断场景。
